% GENERATED FILE. Edit canonical lessons, then run npm run generate:books.
% canonical-source-hash: fnv1a64:751e65e1568c3293
% canonical-lessons: JA-C238-manga, JA-C238-shousetsu, JA-C238-dokusho, JA-C238-bangumi, JA-C238-tsumaranai

\chapter{Comics, Novels and TV}
\label{ch:ja-comics-novels-and-tv}
\hlchaptermodality{238}

\begin{chapteropening}
\textbf{By the end of this chapter:} \emph{I can say what I read and watched for fun.}

Complete the last lesson of chapter 238: I can say what I read and watched for fun.
\end{chapteropening}

\section[\texorpdfstring{\ja{まんが} (manga)}{まんが (manga)}]{\ja{まんが} (manga) — comics; a comic book}

\label{lesson:JA-C238-manga}



\begin{quote}
Before the new one: say the Japanese for to gather, then the Japanese for a doll.
\end{quote}

\subsection*{You'll want to know: \ja{まんが}}
\textbf{\ja{まんが}} — \emph{manga} — \textquotedblleft{}comics; a comic book\textquotedblright{}.

\begin{quote}
\textbf{\ja{きのう} \ja{おとうとは} \ja{まんがを} \ja{よみました。}} — \emph{kinō otōto wa manga o yomimashita} — \textquotedblleft{}Yesterday my younger brother read comics.\textquotedblright{}
\end{quote}

\begin{cousinweb}
\textbf{\ja{まんが}} (\emph{manga}) is Sino-Japanese (kango): two Chinese characters read with their on'yomi, readings borrowed from Chinese. They are \emph{man}, \textquotedblleft{}unrestrained\textquotedblright{} (Middle Chinese \emph{manH}) and \emph{ga}, \textquotedblleft{}picture\textquotedblright{} (\emph{hweaH}) (KANJIDIC2; Middle Chinese in Baxter's transcription).
\end{cousinweb}

\begin{grammarlens}[title={What you've built}]
One more word for everyday life at A2.
\end{grammarlens}

\subsection*{Your turn}
\begin{itemize}
\raggedright
  \item \emph{Say it:} \emph{manga}
  \item \emph{Say it:} \emph{manga}, once more
  \item \emph{Say it:} \emph{kinō otōto wa manga o yomimashita} — Yesterday my younger brother read comics.
  \item \emph{Recall:} say the Japanese for to gather, then the Japanese for a doll, then say \emph{manga} again
\end{itemize}

\begin{tcolorbox}[breakable,colback=teal!4,colframe=teal!35!black,title={Before you move on}]
What does \ja{まんが} mean? (\textquotedblleft{}Comics; a comic book\textquotedblright{}.) Say it once more.
\end{tcolorbox}
\section[\texorpdfstring{\ja{しょうせつ} (shōsetsu)}{しょうせつ (shōsetsu)}]{\ja{しょうせつ} (shōsetsu) — a novel; a story}

\label{lesson:JA-C238-shousetsu}



\begin{quote}
Before the new one: say the Japanese for a doll, then the Japanese for comics.
\end{quote}

\subsection*{You'll want to know: \ja{しょうせつ}}
\textbf{\ja{しょうせつ}} — \emph{shōsetsu} — \textquotedblleft{}a novel; a story\textquotedblright{}.

\begin{quote}
\textbf{\ja{なつやすみに} \ja{ながい} \ja{しょうせつを} \ja{よみました。}} — \emph{natsuyasumi ni nagai shōsetsu o yomimashita} — \textquotedblleft{}In the summer holidays I read a long novel.\textquotedblright{}
\end{quote}

\begin{cousinweb}
\textbf{\ja{しょうせつ}} (\emph{shōsetsu}) is Sino-Japanese (kango): two Chinese characters read with their on'yomi, readings borrowed from Chinese. They are \emph{shō}, \textquotedblleft{}small\textquotedblright{} (Middle Chinese \emph{sjewX}) and \emph{setsu}, \textquotedblleft{}explanation, theory\textquotedblright{} (\emph{sywet}) (KANJIDIC2; Middle Chinese in Baxter's transcription).
\end{cousinweb}

\begin{grammarlens}[title={What you've built}]
One more word for everyday life at A2.
\end{grammarlens}

\subsection*{Your turn}
\begin{itemize}
\raggedright
  \item \emph{Say it:} \emph{shōsetsu}
  \item \emph{Say it:} \emph{shōsetsu}, once more
  \item \emph{Say it:} \emph{natsuyasumi ni nagai shōsetsu o yomimashita} — In the summer holidays I read a long novel.
  \item \emph{Recall:} say the Japanese for a doll, then the Japanese for comics, then say \emph{shōsetsu} again
\end{itemize}

\begin{tcolorbox}[breakable,colback=teal!4,colframe=teal!35!black,title={Before you move on}]
What does \ja{しょうせつ} mean? (\textquotedblleft{}A novel; a story\textquotedblright{}.) Say it once more.
\end{tcolorbox}
\section[\texorpdfstring{\ja{どくしょ} (dokusho)}{どくしょ (dokusho)}]{\ja{どくしょ} (dokusho) — reading; reading books}

\label{lesson:JA-C238-dokusho}



\begin{quote}
Before the new one: say the Japanese for comics, then the Japanese for a novel.
\end{quote}

\subsection*{You'll want to know: \ja{どくしょ}}
\textbf{\ja{どくしょ}} — \emph{dokusho} — \textquotedblleft{}reading; reading books\textquotedblright{}.

\begin{quote}
\textbf{\ja{せんしゅう} \ja{としょかんで} \ja{どくしょを} \ja{しました。}} — \emph{senshū toshokan de dokusho o shimashita} — \textquotedblleft{}Last week I did some reading at the library.\textquotedblright{}
\end{quote}

\begin{cousinweb}
\textbf{\ja{どくしょ}} (\emph{dokusho}) is Sino-Japanese (kango): two Chinese characters read with their on'yomi, readings borrowed from Chinese. They are \emph{doku}, \textquotedblleft{}read\textquotedblright{} (Middle Chinese \emph{duwk}) and \emph{sho}, \textquotedblleft{}write\textquotedblright{} (\emph{syo}) (KANJIDIC2; Middle Chinese in Baxter's transcription).
\end{cousinweb}

\begin{grammarlens}[title={What you've built}]
One more word for everyday life at A2.
\end{grammarlens}

\subsection*{Your turn}
\begin{itemize}
\raggedright
  \item \emph{Say it:} \emph{dokusho}
  \item \emph{Say it:} \emph{dokusho}, once more
  \item \emph{Say it:} \emph{senshū toshokan de dokusho o shimashita} — Last week I did some reading at the library.
  \item \emph{Recall:} say the Japanese for comics, then the Japanese for a novel, then say \emph{dokusho} again
\end{itemize}

\begin{tcolorbox}[breakable,colback=teal!4,colframe=teal!35!black,title={Before you move on}]
What does \ja{どくしょ} mean? (\textquotedblleft{}Reading; reading books\textquotedblright{}.) Say it once more.
\end{tcolorbox}
\section[\texorpdfstring{\ja{ばんぐみ} (bangumi)}{ばんぐみ (bangumi)}]{\ja{ばんぐみ} (bangumi) — a programme; a TV or radio programme}

\label{lesson:JA-C238-bangumi}



\begin{quote}
Before the new one: say the Japanese for a novel, then the Japanese for reading.
\end{quote}

\subsection*{You'll want to know: \ja{ばんぐみ}}
\textbf{\ja{ばんぐみ}} — \emph{bangumi} — \textquotedblleft{}a programme; a TV or radio programme\textquotedblright{}.

\begin{quote}
\textbf{\ja{きのう} \ja{おもしろい} \ja{ばんぐみを} \ja{みました。}} — \emph{kinō omoshiroi bangumi o mimashita} — \textquotedblleft{}Yesterday I watched an interesting programme.\textquotedblright{}
\end{quote}

\begin{cousinweb}
\textbf{\ja{ばんぐみ}} (\emph{bangumi}) is a mixed word: Sino-Japanese \emph{ban}, \textquotedblleft{}turn, number\textquotedblright{} (Middle Chinese \emph{phjon}), borrowed from Chinese, joined to native \emph{kumi}, \textquotedblleft{}set, group\textquotedblright{}, from \emph{kumu}, \textquotedblleft{}to put together\textquotedblright{} (JMdict, KANJIDIC2; Middle Chinese in Baxter's transcription).
\end{cousinweb}

\begin{grammarlens}[title={What you've built}]
One more word for everyday life at A2.
\end{grammarlens}

\subsection*{Your turn}
\begin{itemize}
\raggedright
  \item \emph{Say it:} \emph{bangumi}
  \item \emph{Say it:} \emph{bangumi}, once more
  \item \emph{Say it:} \emph{kinō omoshiroi bangumi o mimashita} — Yesterday I watched an interesting programme.
  \item \emph{Recall:} say the Japanese for a novel, then the Japanese for reading, then say \emph{bangumi} again
\end{itemize}

\begin{tcolorbox}[breakable,colback=teal!4,colframe=teal!35!black,title={Before you move on}]
What does \ja{ばんぐみ} mean? (\textquotedblleft{}A programme; a TV or radio programme\textquotedblright{}.) Say it once more.
\end{tcolorbox}
\section[\texorpdfstring{\ja{つまらない} (tsumaranai)}{つまらない (tsumaranai)}]{\ja{つまらない} (tsumaranai) — boring; dull}

\label{lesson:JA-C238-tsumaranai}



\begin{quote}
Before the new one: say the Japanese for reading, then the Japanese for a programme.
\end{quote}

\subsection*{You'll want to know: \ja{つまらない}}
\textbf{\ja{つまらない}} — \emph{tsumaranai} — \textquotedblleft{}boring; dull\textquotedblright{}.

\begin{quote}
\textbf{\ja{きのう} \ja{つまらない} \ja{えいがを} \ja{みました。}} — \emph{kinō tsumaranai eiga o mimashita} — \textquotedblleft{}Yesterday I watched a boring film.\textquotedblright{}
\end{quote}

\begin{cousinweb}
\textbf{\ja{つまらない}} (\emph{tsumaranai}) is a native Japanese word (wago), inherited rather than borrowed: it is built from \emph{tsumaru}, \textquotedblleft{}to be packed\textquotedblright{}, with the negative ending \emph{-nai}, and it is usually written in kana (JMdict, KANJIDIC2).
\end{cousinweb}

\begin{grammarlens}[title={What you've built}]
One more word for everyday life at A2.
\end{grammarlens}

\subsection*{Your turn}
\begin{itemize}
\raggedright
  \item \emph{Say it:} \emph{tsumaranai}
  \item \emph{Say it:} \emph{tsumaranai}, once more
  \item \emph{Say it:} \emph{kinō tsumaranai eiga o mimashita} — Yesterday I watched a boring film.
  \item \emph{Recall:} say the Japanese for reading, then the Japanese for a programme, then say \emph{tsumaranai} again
\end{itemize}

\begin{tcolorbox}[breakable,colback=teal!4,colframe=teal!35!black,title={Before you move on}]
What does \ja{つまらない} mean? (\textquotedblleft{}Boring; dull\textquotedblright{}.) Say it once more.
\end{tcolorbox}
